\documentclass[12pt,a4paper]{article}

% ================================================================
% PACOTES
% ================================================================

\usepackage[utf8]{inputenc}
\usepackage[T1]{fontenc}
\usepackage[brazilian]{babel}

\usepackage{amsmath}
\usepackage{amssymb}

\usepackage{booktabs}
\usepackage{array}
\usepackage{enumitem}
\usepackage{quoting}
\usepackage{microtype}
\usepackage{geometry}
\usepackage{csquotes}


\usepackage[style=abnt, backend=biber]{biblatex}
\addbibresource{referencias.bib} 

\usepackage{listings}
\usepackage{url}
\usepackage{hyperref}

% ================================================================
% CONFIGURAÇÕES
% ================================================================

\quotingsetup{
    font=small,
    leftmargin=4cm,
    rightmargin=0cm
}

\geometry{
    a4paper,
    margin=2.5cm
}

\hypersetup{
    colorlinks=true,
    linkcolor=black,
    urlcolor=blue,
    citecolor=black
}

% Espaçamento das listas
\setlist[itemize]{
    itemsep=0.2em,
    topsep=0.3em
}

\lstset{
  numbers=left,        % Coloca os números à esquerda
  stepnumber=1,        % Numera de 1 em 1 linha
  numberstyle=\tiny,   % Tamanho da fonte dos números
  numbersep=5pt,       % Espaço entre o número e o código
  literate={ã}{{\~a}}1 {ç}{{\c{c}}}1 {ê}{{\^e}}1 {í}{{\'i}}1
}

% ================================================================
% INFORMAÇÕES DO DOCUMENTO
% ================================================================

\title{
    \textbf{Relatório de Configuração do Ambiente}\\[0.3cm]
    \large Ambiente de Testes de Injeção de Falhas em Dispositivos IoT
}

\author{Eduardo Diniz da Silva}
\date{}

% ================================================================
% DOCUMENTO
% ================================================================

\begin{document}

% ================================================================
% FOLHA DE ROSTO
% ================================================================

\begin{center}
    
    \large
    Escola de Artes, Ciências e Humanidades --- Universidade de São Paulo

    \vspace{1.5cm}

    \Large
    \textbf{RELATÓRIO DE CONFIGURAÇÃO DO AMBIENTE}

    \vspace{0.3cm}

    \large
    \textit{Ambiente de Testes de Injeção de Falhas em Dispositivos IoT}

    \vspace{1cm}

    \normalsize

    \begin{tabular}{ll}
        \textbf{Área:} & Sistemas de Informação / Segurança da Informação \\
        \textbf{Estudante:} & Eduardo Diniz da Silva \\
        \textbf{Orientador:} & Prof. Fábio Nakano
    \end{tabular}

\end{center}

% ================================================================
% RESUMO
% ================================================================

\section{Resumo}

Este relatório descreve a configuração do ambiente de testes de injeção de falhas em dispositivos IoT.


% ================================================================
% CONFIGURAÇÃO INICIAL DO AMBIENTE
% ================================================================

\section{Configuração Inicial do Ambiente}

A primeira etapa foi a instalação do QEMU, emulador de hardware, do GDB, depurador de código, do Newlib, biblioteca C padrão e do Make, ferramenta de automação de compilação.
\begin{lstlisting}[language=bash]
    sudo dnf install qemu-system-arm 
    sudo dnf install gdb 
    sudo dnf install arm-none-eabi-gcc arm-none-eabi-newlib 
    sudo dnf install make
\end{lstlisting} 
Em seguida, foi feito o mapeamento físico da placa virtual no arquivo \texttt{linker.ld}, que define a memória FLASH e RAM, bem como o endereço do topo da pilha (stack) da RAM. Após isso, foi implementado o arquivo 
\texttt{startup.c}, que contém o código de inicialização do sistema. Ambos tiveram como referência o repositório \cite{cortexm_baremetal} e estão disponíveis no repositório do projeto \url{https://github.com/D1n7z/ic-ascon}.
Além disso, foi criado o arquivo \texttt{Makefile}, que automatiza o processo de compilação do projeto:
\begin{itemize}
    \item \texttt{make build}: compila o projeto e gera o binário\texttt{firmware.elf}.
    \item \texttt{make run}: executa o binário no QEMU
    \item \texttt{make clean}: remove os arquivos gerados durante a compilação para uma reconstrução limpa.
\end{itemize}
Finalizando a configuração inicial do ambiente, foi feito um teste de execução do binário no QEMU, através de um arquivo \texttt{main.c} 
\begin{lstlisting} [language=C]
    #include <stdio.h>

    extern void initialise_monitor_handles(void);

    int main(void) {
        initialise_monitor_handles();

        printf("\nHello, World!\n");
        return 0;
    }
\end{lstlisting}
que imprimiu a mensagem "Hello, World!" no console do terminal. Com isso, a parte inicial de instalação e configuração do QEMU foi concluída com sucesso.

% ===============================================================
% CONFIGURAÇÃO DO GDB
% ================================================================

A próxima etapa foi a configuração do GDB, que é o depurador de código. Para isso, foi criado um novo \texttt{main.c} com um código simples de soma para testar o GDB:
\begin{lstlisting} [language=C]
    #include <stdio.h>

    extern void initialise_monitor_handles(void);

    int main(void) {
        initialise_monitor_handles();

        volatile int soma = 0;

        printf("Inicio da execucao do programa\n");
        for (int i = 0; i < 5; i++) {
            soma += i;
        }
        printf("Soma: %d\n", soma);

        return 0;
    }
\end{lstlisting}
O código foi compilado e executado normalmente no QEMU, imprimindo a mensagem "Inicio da execução do programa" e o resultado da soma = 10. Em seguida, foi executado no QEMU com o comando \texttt{qemu-system-arm -M mps2-an385 -cpu cortex-m3 -m 16M -nographic -kernel firmware.elf -semihosting -S -gdb tcp::1234} e o GDB foi iniciado com o comando \texttt{gdb firmware.elf} e conectado ao QEMU com o comando \texttt{target remote :1234}.
Atrvés do GDB, foi possível depurar o código e durante a execução do programa, um registrador teve seu valor alterado de 0 para 100, e o resultado da soma foi alterado de 10 para 110. A fim de validar o ambiente, o processo foi repetido com um script em Python:
\begin{lstlisting}[language=Python]
    import gdb

    gdb.execute("set pagination off")

    print("Injetor de falhas iniciado.")

    # Conectar ao QEMU
    gdb.execute("target remote :1234")

    # Definir breakpoint no início de main
    gdb.execute("break main")
    gdb.execute("continue")

    # Pula para o início do laço for
    for _ in range(3):
        gdb.execute("next")

    # Entrar no laço for
    for _ in range(8):
        gdb.execute("stepi")

    # Lê o valor do registrador r3
    valor_r3 = gdb.parse_and_eval("$r3")
    print(f"Valor original do registrador r3: {valor_r3}")

    # Altera o valor do registrador r3 para 100
    print("Alterando o valor do registrador r3 para 100.")
    gdb.execute("set $r3 = 100")

    # Continua a execução do programa
    print("Continuando a execução do programa")
    gdb.execute("continue")

    # Desconecta do GDB e encerra o script
    gdb.execute("quit")

\end{lstlisting}

O resultado da execução do script foi o mesmo que a execução manual, alterando o valor do registrador r3 de 0 para 100 e o resultado da soma de 10 para 110. Com isso, a configuração do GDB foi concluída com sucesso.


% ================================================================
% REFERÊNCIAS BIBLIOGRÁFICAS
% ================================================================
\printbibliography[title={Referências Bibliográficas}]

\end{document}
